The Optimal Power Flow (OPF) problem is a central tool for the planning and operation of Electric Power Systems (EPS). In the Brazilian context, established industry-grade tools such as CEPEL's ANAREDE incorporate specific control actions \textemdash generator reactive power limits (QLIM), bus voltage magnitude limits (VLIM), automatic switched-shunt susceptance control (CSCA), automatic on-load transformer tap control (CTAP) and distributed energy resource aggregators (DERA) \textemdash that are not natively available in the classical academic formulations implemented by mainstream open-source packages. This work presents the construction, in the Julia language, of an incremental progression of six AC power flow formulations: two reference formulations based on the \textit{PowerModels.jl} package (a full OPF and a classical power flow) and four hand-built \textit{JuMP} formulations in which each of the ANAREDE control actions is sequentially introduced on top of the previous one. Inputs are read from PWF files using the \textit{PWF.jl} package, with extraction of the control structures (DOPC and DLIN) enabled. The non-linear model is solved with the Ipopt solver and the controls are introduced via penalized slack variables in a soft-constraint objective. The six formulations are submitted to a standardized test battery of twenty-five cases, comprising pedagogical networks of three to nine buses, medium-scale networks (300 and 500 buses) and a snapshot of the Brazilian Interconnected Power System (the \textit{01 MAXIMA NOTURNA\_DEZ25} case, with approximately six thousand buses). Results show that the hand-built formulations converge on a strictly larger set of cases than the \textit{PowerModels} references \textemdash in particular, they recover feasible solutions on QLIM-active cases where the classical OPF declares the problem infeasible \textemdash and that the solutions produced by the CSCA- and CTAP-augmented formulations cluster in numerically distinct groups from the reference, reproducing the expected behavior of ANAREDE's control heuristics.

% Separe as Keywords por ponto
\keywords{Optimal power flow. Electric power systems. JuMP. PowerModels.jl. Control actions. Ipopt.}
